我们首先将一般 \problemName 解的结构分解为循环路径的平坦结构，以此开始预备分析。

对于 \problemName 解 $\solution = (\fleetAssignVar, \shareVar, \outsideShareVar)$，我们可以定义其对应的路径方案 $\routingPlan(\solution)$ 或 $\routingPlan(\fleetAssignVar)$ 如下：
\[
    \routingPlan(\solution) = \{\route{f}\mid\forall f\in\fleetSet\}
\]
表示每种资源类型的路径方案，其中 $\route{f}$ 对任意资源类型 $f\in\fleetSet$ 定义如下，表示其路径：
\[
    \route{f} = \{\pathInRoute{f}{k}, \routeAmount{\pathInRoute{f}{k}}\mid k=1,2,\ldots,K\}
\]
\[
    \pathInRoute{f}{k} = [a_1, a_2, \ldots, a_{\lvert \pathInRoute{f}{k} \rvert}]
    \subseteq \{a\in\arcSet\mid \fleetAssignVar_{af}>0\}
    \qquad
    \sum_{k=1}^{\fleetAvail_f}\routeAmount{\pathInRoute{f}{k}}\mathbf{1}(a\in\pathInRoute{f}{k}) = \fleetAssignVar_{af}
    \qquad \forall a\in\arcSet
\]
\[
    \pathInRoute{f}{k}(\servingArc) = \{a\in\servingArc\mid a\in\pathInRoute{f}{k}\}
    ,\qquad
    \pathInRoute{f}{k}(\positionArc) = \{a\in\positionArc\mid a\in\pathInRoute{f}{k}\}
\]
$\pathInRoute{f}{k}$ 表示资源类型 $f$ 的第 $k$ 条路径。$\routeAmount{\pathInRoute{f}{k}}$ 表示 $\pathInRoute{f}{k}$ 在解中出现的次数。注意 $\routeAmount{\pathInRoute{f}{k}}$ 不是指分配到该路径的资源单位数，而是该路径在解中被使用的次数。我们定义 $\wrapnum_{fk}$ 为路径 $\pathInRoute{f}{k}$ 的环绕次数，即该路径跨越规划时域的次数。因此，分配到路径 $\pathInRoute{f}{k}$ 的资源单位数为 $\routeAmount{\pathInRoute{f}{k}} \cdot \wrapnum_{fk}$。

% ========================

然后，我们建立聚合映射 $\Phi$ 的核心性质。
\paragraph*{$\Phi_{\arcSet}$ 的性质。}
因为 $\Phi_{\arcSet}$ 是一个函数，族 $\{\Phi_{\arcSet}^{-1}(a)\}_{a\in\arcSet}\cup\{\Phi_{\arcSet}^{-1}(\bot)\}$ 构成 $\full{\arcSet}$ 的一个划分。
此外，每条服务弧 $a'\in\full{\flightArc}$ 映射到 $\flightArc$ 的唯一元素。因此
\[
    \full{\flightArc} = \biguplus_{a\in\flightArc}\Phi_{\arcSet}^{-1}(a).
\]
对于停留弧，我们得到不相交分解
\[
    \full{\holdingArc} = \left(\biguplus_{a\in\holdingArc}\Phi_{\arcSet}^{-1}(a)\right) \uplus \Phi_{\arcSet}^{-1}(\bot),
\]
其中 $\Phi_{\arcSet}^{-1}(\bot)$ 包含完全被吸收在单个部分时间节点处的短暂地面等待。
节点映射 $\Phi_{\nodeSet}$ 具有相同的划分性质：$\full{\nodeSet} = \biguplus_{(i,t)\in\nodeSet} \Phi_{\nodeSet}^{-1}(i,t)$。

\paragraph*{原像的相容性。}
固定 $a=((i,t)\rightarrow(j,t'))\in\arcSet$，设 $a'=((i,\tau)\rightarrow(j,\tau'))$ 为完整网络中的一条弧。
由 $\Phi_{\arcSet}$ 的定义，我们有
\[
    a'\in\Phi_{\arcSet}^{-1}(a)
    \iff 
    \Phi_{\nodeSet}((i,\tau))=(i,t)\ \text{and}\ \Phi_{\nodeSet}((j,\tau'))=(j,t').
\]
因此 $a'$ 的尾部自动位于 $\Phi_{\nodeSet}^{-1}(i,t)$ 中，其头部位于 $\Phi_{\nodeSet}^{-1}(j,t')$ 中，反之，任何从 $(i,\tau)\in\Phi_{\nodeSet}^{-1}(i,t)$ 出发到达 $(j,\tau')\in\Phi_{\nodeSet}^{-1}(j,t')$ 的弧必须映射到 $a$（而非 $\bot$）。
等价地，
\[
    a\in\outArcSet{(i,t)},\ a'\in\Phi_{\arcSet}^{-1}(a)
    \iff \exists (i,\tau)\in\Phi_{\nodeSet}^{-1}(i,t)\ \text{s.t.}\ a'\in\outArcSet{(i,\tau)},
\]
and similarly
\[
    a\in\inArcSet{(j,t')},\ a'\in\Phi_{\arcSet}^{-1}(a)
    \iff \exists (j,\tau')\in\Phi_{\nodeSet}^{-1}(j,t')\ \text{s.t.}\ a'\in\inArcSet{(j,\tau')}.
\]
这种相容性使我们能够在部分弧与其完整网络原像之间交换求和。

特别地，由于原像族 $\{\Phi_{\arcSet}^{-1}(a)\}_{a\in\flightArc}$ 构成完整服务弧集 $\full{\flightArc}$ 的严格数学划分，全求和律成立。对于以完整服务弧 $a'\in\full{\flightArc}$ 为索引的任意实值参数或变量 $y_{a'}$，我们有：
\[
    \sum_{a\in\flightArc}\ \sum_{a'\in\Phi_{\arcSet}^{-1}(a)} y_{a'} 
    = \sum_{a'\in\full{\flightArc}} y_{a'}.
\]

预备工作的最后一个组成部分是路由存在性引理，它对于将宏观流量分解为具体循环至关重要。
\begin{lemma}
    对于任意 $\model{\discretization}$-可行解 $\solution$，总可以找到对应的路由方案 $\routingPlan(\solution)$。

\end{lemma}

\begin{proof}
    \newcommand{\assignmentVector}{\mathbf{x}}
    \newcommand{\feasibleSet}{\mathscr{X}}
    \newcommand{\vectorAdd}{\oplus }
    \newcommand{\vectorMinus}{\ominus}
    \newcommand{\vectorLe}{\preceq}
We fix a resource type $f\in\fleetSet$ and prove the statement for its assignment vector 
$\assignmentVector^f=\{\fleetAssignVar_{af}\}_{a\in\arcSet}$. $\assignmentVector^f_{a}$ denotes the number of resource type $f$ assigned to arc $a$, which is a non-negative real number.

\paragraph{Step 0. Algebraic structure.}
Define the flow feasible set denoting all assignment vectors satisfying the flow-balance constraint~\eqref{constr:main.flow}:
\[
\feasibleSet_f=\Bigl\{
\assignmentVector^f\in\mathbb{R}_{\ge0}^{\lvert\arcSet\rvert}\;\Bigm|\;
\sum_{a\in\outArcSet{i,t}}\assignmentVector^f_{a}-\sum_{a\in\inArcSet{i,t}}\assignmentVector^f_{a}=0,\ \forall(i,t)\in\nodeSet
\Bigr\}.
\]
\[
    \vectorAdd: \feasibleSet_f \times \feasibleSet_f \to \feasibleSet_f,\quad x \vectorAdd x' = \{x_a + x'_a\}_{a\in\arcSet}
\]
\[
    \vectorMinus: \feasibleSet_f \times \feasibleSet_f \to \feasibleSet_f,\quad x \vectorMinus x' = \{x_a - x'_a\}_{a\in\arcSet}
\]
\[
    \vectorLe: x \vectorLe x' \iff x_a \le x'_a,\ \forall a\in\arcSet
\]
$(\feasibleSet_f,\vectorAdd)$ forms a commutative monoid:
it is closed under addition $\vectorAdd$, associative, and admits the zero vector $\mathbf{0}$ as identity.
We equip it with the partial order $\vectorLe$ defined componentwise; for $x',x\in\feasibleSet_f$,
the difference $x\vectorMinus x'$ is defined if $x'\vectorLe x$ and remains in $\feasibleSet_f$.

\paragraph{Step 1. Cycle vectors.}
A \emph{directed cycle} $C=(a_1,\ldots,a_L)$ is a finite sequence of arcs 
forming a closed walk on $\net$; the same arc may appear multiple times along $C$.
For any arc $a\in\arcSet$, denote by
\[
    \nu_{C}(a) := \#\{r\in\{1,\ldots,L\}\mid a_r = a\}
\]
its multiplicity within $C$.
For any scalar $z\ge 0$, define its \emph{cycle vector}
\[
v(C;z)_a=
\begin{cases}
z\,\nu_{C}(a),& a\in C,\\[2pt]
0,&\text{otherwise}.
\end{cases}
\]
Then $v(C;z)\in\feasibleSet_f$ because every node on the cycle has one incoming and one outgoing arc,
and scaling by $z$ preserves the flow-balance constraint~\eqref{constr:main.flow}.

\paragraph{Step 2. Existence of a directed cycle.}
For a nonzero $\assignmentVector\in\feasibleSet_f$, define the positive-support subgraph 
${\net}^+(\assignmentVector)$ induced by $\arcSet^+(\assignmentVector)=\{a\in\arcSet\mid \assignmentVector_a>0\}$.
At every node of ${\net}^+(\assignmentVector)$, the flow-balance constraint ensures
$\text{outdegree}\ge1$ and $\text{indegree}\ge1$.
Starting from any node with positive outflow and following positive arcs,
finiteness of the network implies that a node must repeat, 
thus forming a directed cycle $C\subseteq \arcSet^+(\assignmentVector)$.

\paragraph{Step 3. Greedy peeling procedure.}
Let $z=\min_{a\in C}\assignmentVector_a/\nu_{C}(a)>0$ (so that $v(C;z)\vectorLe \assignmentVector$) and define $\assignmentVector':=\assignmentVector-v(C;z)$.
Then:
\begin{itemize}
\item $\assignmentVector'\ge0$ componentwise, because $z$ is the minimum scaled value $\assignmentVector_a/\nu_C(a)$;
\item $\assignmentVector'\in\feasibleSet_f$, since both $\assignmentVector$ and $v(C;z)$ satisfy~\eqref{constr:main.flow};
\item the support strictly shrinks: $\arcSet^+(\assignmentVector')\subsetneq \arcSet^+(\assignmentVector)$.
\end{itemize}
Therefore, for $\assignmentVector^f$ iterating this procedure finitely many times (at most $|A^+(\assignmentVector^f)|$)
yields a decomposition
\[
\assignmentVector^f=\sum_{k=1}^{K}v(C_k;z_k),\qquad K\le|A^+(\assignmentVector^f)|.
\]

\paragraph{Step 4. Constructing the routing plan.}
For each cycle $C_k$, let $\pathInRoute{f}{k}$ be its projection onto the set of service arcs
$\flightArc$.
Let $\routeAmount{\pathInRoute{f}{k}}:=z_k$, and denote
\[
    \nu_{fk}(a)
        := \#\{r\mid a_{k,r}=a\}
        = \nu_{C_k}(a),
\]
the number of times arc $a$ appears in the projected cycle (if $a\notin C_k$, the count is zero).
Then, for any $a\in\flightArc$,
\[
\sum_{k=1}^{K}\routeAmount{\pathInRoute{f}{k}}\nu_{fk}(a)
=\sum_{k=1}^{K}v(C_k;z_k)_a
=\assignmentVector^f_a
=\fleetAssignVar_{af},
\]
Define 
$\route{f}=\{(\pathInRoute{f}{k},\routeAmount{\pathInRoute{f}{k}})\}_{k=1}^{K}$
and 
$\routingPlan(\solution)=\{\route{f}\mid f\in\fleetSet\}$.
Hence, for every resource type $f$, we can derive a routing plan consistent with its assignment vector.
\end{proof}

接下来，我们引入路径存在性引理，该引理保证任何流量平衡的资源类型分配都可以分解为具体的循环路径。
\begin{lemma}
    对于任意 $\model{\discretization}$-可行解 $\solution$，总可以找到对应的路由方案 $\routingPlan(\solution)$。

\end{lemma}

\begin{proof}
    \newcommand{\assignmentVector}{\mathbf{x}}
    \newcommand{\feasibleSet}{\mathscr{X}}
    \newcommand{\vectorAdd}{\oplus }
    \newcommand{\vectorMinus}{\ominus}
    \newcommand{\vectorLe}{\preceq}
We fix a resource type $f\in\fleetSet$ and prove the statement for its assignment vector 
$\assignmentVector^f=\{\fleetAssignVar_{af}\}_{a\in\arcSet}$. $\assignmentVector^f_{a}$ denotes the number of resource type $f$ assigned to arc $a$, which is a non-negative real number.

\paragraph{Step 0. Algebraic structure.}
Define the flow feasible set denoting all assignment vectors satisfying the flow-balance constraint~\eqref{constr:main.flow}:
\[
\feasibleSet_f=\Bigl\{
\assignmentVector^f\in\mathbb{R}_{\ge0}^{\lvert\arcSet\rvert}\;\Bigm|\;
\sum_{a\in\outArcSet{i,t}}\assignmentVector^f_{a}-\sum_{a\in\inArcSet{i,t}}\assignmentVector^f_{a}=0,\ \forall(i,t)\in\nodeSet
\Bigr\}.
\]
\[
    \vectorAdd: \feasibleSet_f \times \feasibleSet_f \to \feasibleSet_f,\quad x \vectorAdd x' = \{x_a + x'_a\}_{a\in\arcSet}
\]
\[
    \vectorMinus: \feasibleSet_f \times \feasibleSet_f \to \feasibleSet_f,\quad x \vectorMinus x' = \{x_a - x'_a\}_{a\in\arcSet}
\]
\[
    \vectorLe: x \vectorLe x' \iff x_a \le x'_a,\ \forall a\in\arcSet
\]
$(\feasibleSet_f,\vectorAdd)$ forms a commutative monoid:
it is closed under addition $\vectorAdd$, associative, and admits the zero vector $\mathbf{0}$ as identity.
We equip it with the partial order $\vectorLe$ defined componentwise; for $x',x\in\feasibleSet_f$,
the difference $x\vectorMinus x'$ is defined if $x'\vectorLe x$ and remains in $\feasibleSet_f$.

\paragraph{Step 1. Cycle vectors.}
A \emph{directed cycle} $C=(a_1,\ldots,a_L)$ is a finite sequence of arcs 
forming a closed walk on $\net$; the same arc may appear multiple times along $C$.
For any arc $a\in\arcSet$, denote by
\[
    \nu_{C}(a) := \#\{r\in\{1,\ldots,L\}\mid a_r = a\}
\]
its multiplicity within $C$.
For any scalar $z\ge 0$, define its \emph{cycle vector}
\[
v(C;z)_a=
\begin{cases}
z\,\nu_{C}(a),& a\in C,\\[2pt]
0,&\text{otherwise}.
\end{cases}
\]
Then $v(C;z)\in\feasibleSet_f$ because every node on the cycle has one incoming and one outgoing arc,
and scaling by $z$ preserves the flow-balance constraint~\eqref{constr:main.flow}.

\paragraph{Step 2. Existence of a directed cycle.}
For a nonzero $\assignmentVector\in\feasibleSet_f$, define the positive-support subgraph 
${\net}^+(\assignmentVector)$ induced by $\arcSet^+(\assignmentVector)=\{a\in\arcSet\mid \assignmentVector_a>0\}$.
At every node of ${\net}^+(\assignmentVector)$, the flow-balance constraint ensures
$\text{outdegree}\ge1$ and $\text{indegree}\ge1$.
Starting from any node with positive outflow and following positive arcs,
finiteness of the network implies that a node must repeat, 
thus forming a directed cycle $C\subseteq \arcSet^+(\assignmentVector)$.

\paragraph{Step 3. Greedy peeling procedure.}
Let $z=\min_{a\in C}\assignmentVector_a/\nu_{C}(a)>0$ (so that $v(C;z)\vectorLe \assignmentVector$) and define $\assignmentVector':=\assignmentVector-v(C;z)$.
Then:
\begin{itemize}
\item $\assignmentVector'\ge0$ componentwise, because $z$ is the minimum scaled value $\assignmentVector_a/\nu_C(a)$;
\item $\assignmentVector'\in\feasibleSet_f$, since both $\assignmentVector$ and $v(C;z)$ satisfy~\eqref{constr:main.flow};
\item the support strictly shrinks: $\arcSet^+(\assignmentVector')\subsetneq \arcSet^+(\assignmentVector)$.
\end{itemize}
Therefore, for $\assignmentVector^f$ iterating this procedure finitely many times (at most $|A^+(\assignmentVector^f)|$)
yields a decomposition
\[
\assignmentVector^f=\sum_{k=1}^{K}v(C_k;z_k),\qquad K\le|A^+(\assignmentVector^f)|.
\]

\paragraph{Step 4. Constructing the routing plan.}
For each cycle $C_k$, let $\pathInRoute{f}{k}$ be its projection onto the set of service arcs
$\flightArc$.
Let $\routeAmount{\pathInRoute{f}{k}}:=z_k$, and denote
\[
    \nu_{fk}(a)
        := \#\{r\mid a_{k,r}=a\}
        = \nu_{C_k}(a),
\]
the number of times arc $a$ appears in the projected cycle (if $a\notin C_k$, the count is zero).
Then, for any $a\in\flightArc$,
\[
\sum_{k=1}^{K}\routeAmount{\pathInRoute{f}{k}}\nu_{fk}(a)
=\sum_{k=1}^{K}v(C_k;z_k)_a
=\assignmentVector^f_a
=\fleetAssignVar_{af},
\]
Define 
$\route{f}=\{(\pathInRoute{f}{k},\routeAmount{\pathInRoute{f}{k}})\}_{k=1}^{K}$
and 
$\routingPlan(\solution)=\{\route{f}\mid f\in\fleetSet\}$.
Hence, for every resource type $f$, we can derive a routing plan consistent with its assignment vector.
\end{proof}

路径分解建立后，我们定义若干辅助关系和集合，以便于后续对跨时间截面的资源计数进行评估。

\subparagraph*{由 $\Phi_{\arcSet}$ 导出的原像记号。}
设 $\Phi_{\arcSet}:\full{\arcSet}\to\widetilde{\arcSet}$ 为定义松弛解时构造的弧映射。
对 $\forall \mathcal{B}\subseteq\arcSet$ 记
\[
    \Phi_{\arcSet}^{-1}(\mathcal{B})
        := \biguplus_{a\in\mathcal{B}}\Phi_{\arcSet}^{-1}(a)
        \subseteq \full{\arcSet}.
\]
显然，
\[
    a'\in\Phi_{\arcSet}^{-1}(\mathcal{B})
        \iff \Phi_{\arcSet}(a')\in\mathcal{B}.
\]

\subparagraph*{辅助集合。}
预期从完整网络解 $\bar{\fleetAssignVar}$ 到松弛解 $\fleetAssignVar$ 的映射，考虑聚合公式：
\[
    \fleetAssignVar_{af}
        = \sum_{a'\in\Phi_{\arcSet}^{-1}(a)}\bar{\fleetAssignVar}_{a'f},
    \qquad
    x_{af}\ge0,
    \label{eq:fleet_avail.aggregate}
\]
由上述路径存在性引理，完整网络流 $\bar{\fleetAssignVar}_{\cdot f}$ 允许路径分解：
\[
    \bar{\fleetAssignVar}_{a'f}
        = \sum_{k\in K_f}
          \routeAmount{\pathInRoute{f}{k}}\,
          \nu_{fk}(a'),
    \label{eq:fleet_avail.route_decomp}
\]
其中 $\pathInRoute{f}{k}=(a'_{k,1},\ldots,a'_{k,q_k})$ 表示资源类型 $f$ 按时间顺序排列的第 $k$ 条闭合游走，且 $\nu_{fk}(a')=\#\{r\mid a'_{k,r}=a'\}$。
逐分量应用 $\Phi_{\arcSet}$ 得到部分弧的循环序列 $\widetilde{\pathInRoute{f}{k}}=(a_{k,1},\ldots,a_{k,Q_k})$，其中当像不为 $\bot$ 时 $a_{k,r}=\Phi_{\arcSet}(a'_{k,r})$。
对每个 $k$ 定义：
\begin{equation}
\label{eq:aux_sets_def}
\begin{aligned}
    \wrapSet_{fk}
        &:= \left\{a_{k,r}\in\widetilde{\pathInRoute{f}{k}}
            \mid t_o(a_{k,r})+\travelTime_{a_{k,r}}\ge \planningEnd\right\},\\
    \arcPosi
        &:= \left\{a_{k,r}\in\widetilde{\pathInRoute{f}{k}}
            \mid t_o(a_{k,r})\preceq_{t_o(a_{k,r})} \checkTime
                 \prec_{t_o(a_{k,r})} t_d(a_{k,r})\right\},\\
    \arcNega
        &:= \left\{a_{k,r}\in\widetilde{\pathInRoute{f}{k}}
            \mid t_d(a_{k,r})\preceq_{t_d(a_{k,r})} \checkTime
                 \prec_{t_d(a_{k,r})} t_o(a_{k,r}),\;
                 t_o(a_{k,r})+\travelTime_{a_{k,r}} < \planningEnd\right\},
\end{aligned}
\end{equation}
并记 $\wrapnum_{fk}=|\wrapSet_{fk}|$。
设 $H:=\planningEnd-\planningStart$ 为时域长度。
对每个基准点 $u$，在 $[0,H)$ 上定义循环关系如下：
\begin{equation*}
\begin{aligned}
    \text{（更早）}\qquad t_1\prec_u t_2 &\iff t_1-u < t_2-u \pmod{H},\\
    \text{（不晚于）}\qquad t_1\preceq_u t_2 &\iff t_1-u \le t_2-u \pmod{H},\\
    \text{（更晚）}\qquad t_1\succ_u t_2 &\iff t_1-u > t_2-u \pmod{H},\\
    \text{（不早于）}\qquad t_1\succeq_u t_2 &\iff t_1-u \ge t_2-u \pmod{H}.
\end{aligned}
\end{equation*}


\begin{lemma}\label{lem:fleet_avail.expand}
For every resource type $f\in\fleetSet$ and cut time $\checkTime$,
\begin{equation*}
\begin{aligned}
    &\sum_{a\in\holdingArcTc{\checkTime}}\fleetAssignVar_{af}
     + \sum_{a\in\flightArcTc{\checkTime}}\fleetAssignVar_{af}
     - \sum_{a\in\invertArcTc{\checkTime}}\fleetAssignVar_{af}\\
    &=
    \sum_{k\in K_f}
        \routeAmount{\pathInRoute{f}{k}}
        \left(
            |\arcPosi|-|\arcNega|
        \right).
\end{aligned}
\end{equation*}
\end{lemma}
\begin{proof}
Using \eqref{eq:fleet_avail.aggregate} and the disjointness of the family
$\{\Phi_{\arcSet}^{-1}(a)\}_{a\in\arcSet}$ we obtain
\begin{equation*}
\begin{aligned}
    &\sum_{a\in\holdingArcTc{\checkTime}}\fleetAssignVar_{af}
     + \sum_{a\in\flightArcTc{\checkTime}}\fleetAssignVar_{af}
     - \sum_{a\in\invertArcTc{\checkTime}}\fleetAssignVar_{af}\\
    &=
    \sum_{a'\in\Phi_{\arcSet}^{-1}(\holdingArcTc{\checkTime})}\bar{\fleetAssignVar}_{a'f}
    +\sum_{a'\in\Phi_{\arcSet}^{-1}(\flightArcTc{\checkTime})}\bar{\fleetAssignVar}_{a'f}
    -\sum_{a'\in\Phi_{\arcSet}^{-1}(\invertArcTc{\checkTime})}\bar{\fleetAssignVar}_{a'f}.
\end{aligned}
\end{equation*}
Inserting the route decomposition~\eqref{eq:fleet_avail.route_decomp} and exchanging summations yields
\begin{equation*}
\begin{aligned}
    &\sum_{a\in\holdingArcTc{\checkTime}}\fleetAssignVar_{af}
     + \sum_{a\in\flightArcTc{\checkTime}}\fleetAssignVar_{af}
     - \sum_{a\in\invertArcTc{\checkTime}}\fleetAssignVar_{af}\\
    % &=          \sum_{a'\in\Phi_{\arcSet}^{-1}(\holdingArcTc{\checkTime})}\sum_{k\in K_f}\routeAmount{\pathInRoute{f}{k}}\#\{r \mid a^{\prime}_{k,r}=a'\}\\
    %     &\quad +\sum_{a'\in\Phi_{\arcSet}^{-1}(\flightArcTc{\checkTime})}\sum_{k\in K_f}\routeAmount{\pathInRoute{f}{k}}\#\{r \mid a^{\prime}_{k,r}=a'\}\\
    %     &\quad -\sum_{a'\in\Phi_{\arcSet}^{-1}(\invertArcTc{\checkTime})}\sum_{k\in K_f}\routeAmount{\pathInRoute{f}{k}}\#\{r \mid a^{\prime}_{k,r}=a'\}\\
    &=
    \sum_{k\in K_f}\routeAmount{\pathInRoute{f}{k}}
        \left(
            \sum_{a'\in\Phi_{\arcSet}^{-1}(\holdingArcTc{\checkTime})}\#\{r \mid a^{\prime}_{k,r}=a'\}
            +\sum_{a'\in\Phi_{\arcSet}^{-1}(\flightArcTc{\checkTime})}\#\{r \mid a^{\prime}_{k,r}=a'\}
            -\sum_{a'\in\Phi_{\arcSet}^{-1}(\invertArcTc{\checkTime})}\#\{r \mid a^{\prime}_{k,r}=a'\}
        \right)\\
    &=
    \sum_{k\in K_f}\routeAmount{\pathInRoute{f}{k}}
        \left(
            \#\{r \mid a^{\prime}_{k,r}\in\Phi_{\arcSet}^{-1}(\holdingArcTc{\checkTime})\}
            +\#\{r \mid a^{\prime}_{k,r}\in\Phi_{\arcSet}^{-1}(\flightArcTc{\checkTime})\}
            -\#\{r \mid a^{\prime}_{k,r}\in\Phi_{\arcSet}^{-1}(\invertArcTc{\checkTime})\}
        \right)\\
    &=
    \sum_{k\in K_f}
        \routeAmount{\pathInRoute{f}{k}}
        \left(
            \#\{r \mid a_{k,r}\in\holdingArcTc{\checkTime}\}
            + \#\{r \mid a_{k,r}\in\flightArcTc{\checkTime}\}
            - \#\{r \mid a_{k,r}\in\invertArcTc{\checkTime}\}
        \right).
\end{aligned}
\end{equation*}
For notational convenience lift every arrival time across the planning horizon:
\[
    \widetilde{t}_o(a) := t_o(a),\qquad
    \widetilde{t}_d(a) :=
        \begin{cases}
            t_d(a),      & t_o(a)+\travelTime_a < \planningEnd,\\
            t_d(a)+H, & t_o(a)+\travelTime_a \ge \planningEnd.
        \end{cases}
\]
Then the circular relations used in the definitions of $\arcPosi$ and $\arcNega$ rewrite as
\begin{equation*}
\begin{aligned}
    a\in\arcPosi
        &\iff \widetilde{t}_o(a) \le \checkTime < \widetilde{t}_d(a),\\
    a\in\arcNega
        &\iff \widetilde{t}_d(a) \le \checkTime < \widetilde{t}_o(a) \wedge t_o(a)+\travelTime_a < \planningEnd.
\end{aligned}
\end{equation*}
Fix $r$ and abbreviate $a:=a_{k,r}$. We compare the lifted inequalities with the constraint sets:
\begin{itemize}
    \item If $a\in\flightArcTc{\checkTime}$, then either
    $t_o(a)\le \checkTime < t_d(a)$ (non wrap) or
    $t_d(a)<t_o(a)\le \checkTime < \planningEnd \le t_o(a)+\travelTime_a$ (wrap case).
    In both situations $\widetilde{t}_o(a) \le \checkTime < \widetilde{t}_d(a)$, hence $a\in\arcPosi$.
    The same implication holds for holding arcs.
    Conversely, if $a\in\arcPosi$ and $t_o(a)+\travelTime_a < \planningEnd$, then $\checkTime$ lies between $t_o(a)$ and $t_d(a)$ without wrapping, so $a$ belongs to $\holdingArcTc{\checkTime}\cup\flightArcTc{\checkTime}$.
    If $a\in\arcPosi$ and $t_o(a)+\travelTime_a \ge \planningEnd$, then $\checkTime$ sits between $\widetilde{t}_o(a)=t_o(a)$ and $\widetilde{t}_d(a)=t_d(a)+H$, forcing the wrap-around condition in $\holdingArcTc{\checkTime}$ or $\flightArcTc{\checkTime}$. Thus
    \[
        a\in\holdingArcTc{\checkTime}\cup\flightArcTc{\checkTime}
        \iff a\in\arcPosi.
    \]
    \item If $a\in\invertArcTc{\checkTime}$, Definition~\ref{def:relax.invertArcTc} furnishes
    $t_d(a)\le \checkTime < t_o(a)$ together with $t_o(a)+\travelTime_a<\planningEnd$.
    Hence $\widetilde{t}_d(a)=t_d(a)$, yielding $\widetilde{t}_d(a) \le \checkTime < \widetilde{t}_o(a)$, i.e.\ $a\in\arcNega$.
    Conversely, if $a\in\arcNega$ then by definition $t_o(a)+\travelTime_a<\planningEnd$ and
    $\widetilde{t}_d(a) \le \checkTime < \widetilde{t}_o(a)$.
    Because the wrap condition excludes $\widetilde{t}_d(a)=t_d(a)+H$, the lifted inequality collapses to $t_d(a)\le \checkTime < t_o(a)$, so $a\in\invertArcTc{\checkTime}$.
\end{itemize}
Therefore
\[
    \#\{r \mid a_{k,r}\in\holdingArcTc{\checkTime}\cup\flightArcTc{\checkTime}\}=|\arcPosi|,
    \qquad
    \#\{r \mid a_{k,r}\in\invertArcTc{\checkTime}\}=|\arcNega|,
\]
and the parenthesis equals $|\arcPosi|-|\arcNega|$.
This proves the lemma.
\end{proof}

\begin{lemma}\label{lem:fleet_avail.wrap}
For every resource type $f$ and route $\pathInRoute{f}{k}$,
$|\arcPosi|-|\arcNega|=\wrapnum_{fk}$.
\end{lemma}
\begin{proof}

We only consider a fixed $k\in K_f$ in the following proof and omit the subscript $k$ for brevity. Denote $\widetilde{\pathInRoute{f}{k}}$ as $(a_{1},a_{2},\ldots,a_{Q})$. Denote $\wrapSet = \wrapSet_{fk}$.
Define operations $\oplus,\,\ominus$ for indexes as follows:
\[
    a\oplus b := a+b \mod Q,
    \qquad
    a\ominus b := a-b \mod Q.
\]

Define the sign profile recording the position of the departure time of arc $a_{r}$ relative to the cut at $\checkTime$:
\[
    \sign_{r} :=
    \begin{cases}
        0, & t_o(a_{r})\le\checkTime,\\
        1, & \text{otherwise}
    \end{cases}
    \qquad r=1,\ldots,Q.
\]

Because the route is continuous, $t_o(a_{r})$ and $t_d(a_{r\ominus1})$ must be equal at the cut. Consider the combination of two consecutive signs, by definitions above, define:
\[
    \slr{r} := (\sign_{r\ominus1},\sign_{r}) = (0,1)\qquad( \iff t_o(a_r)\le\checkTime < t_d(a_r))
\]
\[
    \srl{r} := (\sign_{r\ominus1},\sign_{r}) = (1,0)\qquad( \iff t_d(a_r)\le\checkTime < t_o(a_r))
\]
% \[
%     \sll{r} := (\sign_{r\ominus1},\sign_{r}) = (0,0) \iff t_o(a_r)\le\checkTime \cap t_d(a_r)\le\checkTime
% \]
% \[
%     \srr{r} := (\sign_{r\ominus1},\sign_{r}) = (1,1) \iff t_o(a_r)>\checkTime \cap t_d(a_r)>\checkTime
% \]
\[
    a_r\in\arcPosi \iff \slr{r} \vee (\srr{r} \wedge a_r\in\wrapSet)
\]
\[
    a_r\in\wrapSet \iff (\srr{r} \wedge a_r\in\arcPosi) \vee (\srl{r} \wedge \neg a_r\in\arcNega)
\]
\[
    a_r\in\arcNega \iff \srl{r} \wedge \neg(a_r\in\wrapSet)
\]

Define set $\wrapTc = \arcPosi \cap \wrapSet$. Therefore, we can derive the following equations.
First, we derive that $a_r\in\arcPosi\setminus\wrapTc$ is equivalent to $\slr{r}$:
\begin{equation*}
\begin{aligned}
    a_r\in\arcPosi\setminus\wrapTc
    & \iff a_r\in\arcPosi\setminus\wrapSet \\
    & \iff a_r\in\arcPosi - a_r\in\wrapSet \\
    & \iff a_r\in\arcPosi - ((\srr{r} \wedge a_r\in\arcPosi) \vee (\srl{r} \wedge a_r\in\arcNega)) \\
    & \iff
    (a_r\in\arcPosi - (\srr{r} \wedge a_r\in\arcPosi))
    % \underbrace{
    %     (a_r\in\arcPosi - (\srr{r} \wedge a_r\in\arcPosi))
    % }_{
    %     \boxed{\begin{aligned}
    %         &\iff a_r\in\arcPosi \wedge (\neg\srr{r} \vee \neg a_r\in\arcPosi) \\
    %         &\iff a_r\in\arcPosi \wedge \neg\srr{r} \\
    %         &\iff (\slr{r} \vee (\srr{r} \wedge a_r\in\wrapSet)) \wedge \neg\srr{r} \\
    %         &\iff (\slr{r} \wedge \neg\srr{r}) \vee \mathbf{0} \\
    %         &\iff \slr{r}
    %     \end{aligned}}
    % }
    \wedge
    (a_r\in\arcPosi - (\srl{r} \wedge a_r\in\arcNega))
    % \underbrace{
    %     (a_r\in\arcPosi - (\srl{r} \wedge a_r\in\arcNega))
    % }_{
    %     \boxed{\begin{aligned}
    %         &\iff a_r\in\arcPosi \wedge (\neg\srl{r} \vee \neg a_r\in\arcNega) \\
    %         &\iff (a_r\in\arcPosi \wedge \neg\srl{r}) \vee \mathbf{0} \\
    %         &\iff a_r\in\arcPosi
    %     \end{aligned}}
    % }
    \\
    &
    % \iff
    \overset{\text{(*)}}{\iff}
    \slr{r} \wedge a_r\in\arcPosi \\
    & \iff \slr{r}
\end{aligned}
\end{equation*}
The equivalence marked with $(*)$ holds because:
\begin{equation*}
\begin{aligned}
    a_r\in\arcPosi - (\srr{r} \wedge a_r\in\arcPosi) 
    &\iff a_r\in\arcPosi \wedge (\neg\srr{r} \vee \neg a_r\in\arcPosi) \\
    &\iff a_r\in\arcPosi \wedge \neg\srr{r} \\
    &\iff (\slr{r} \vee (\srr{r} \wedge a_r\in\wrapSet)) \wedge \neg\srr{r} \\
    &\iff (\slr{r} \wedge \neg\srr{r}) \vee \mathbf{0} \\
    &\iff \slr{r}
\end{aligned}
\end{equation*}
and
\begin{equation*}
\begin{aligned}
    a_r\in\arcPosi - (\srl{r} \wedge a_r\in\arcNega)
    &\iff a_r\in\arcPosi \wedge (\neg\srl{r} \vee \neg a_r\in\arcNega) \\
    &\iff (a_r\in\arcPosi \wedge \neg\srl{r}) \vee \mathbf{0} \\
    &\iff a_r\in\arcPosi
\end{aligned}
\end{equation*}

Second, we derive that $a_r\in(\wrapSet\setminus\wrapTc)\cup\arcNega$ is equivalent to $\srl{r}$:
\begin{equation*}
\begin{aligned}
    a_r\in\wrapSet\setminus\wrapTc
    & \iff a_r\in\wrapSet\setminus\arcPosi \\
    & \iff a_r\in\wrapSet - a_r\in\arcPosi \\
    & \iff ((\srr{r} \wedge a_r\in\arcPosi) \vee (\srl{r} \wedge \neg a_r\in\arcNega)) - a_r\in\arcPosi \\
    & \iff \underbrace{((\srr{r} \wedge a_r\in\arcPosi) - a_r\in\arcPosi)}_{\mathbf{0}} \vee ((\srl{r} \wedge \neg a_r\in\arcNega) - a_r\in\arcPosi) \\
    & \iff \srl{r} \wedge \neg a_r\in\arcNega \wedge \neg a_r\in\arcPosi \\
    & \iff \srl{r} \wedge \neg a_r\in\arcNega
\end{aligned}
\end{equation*}

Thus,
\begin{equation*}
\begin{aligned}
    a_r\in(\wrapSet\setminus\wrapTc)\cup\arcNega
    & \iff (\srl{r} \wedge \neg a_r\in\arcNega) \vee a_r\in\arcNega \\
    & \iff \srl{r} \quad (\because a_r\in\arcNega \implies \srl{r})
\end{aligned}
\end{equation*}

Now we have two important results:
\[
    a_r\in\arcPosi\setminus\wrapTc \iff \slr{r}
\]
\[
    a_r\in(\wrapSet\setminus\wrapTc)\cup\arcNega \iff \srl{r}
\]

Because the route is cyclic and continuous, the number of $\slr{r}$ equals the number of $\srl{r}$ when scanning through the sequence:
\[
    \#\{r\mid \slr{r}\} = \#\{r\mid \srl{r}\}
\]

which means $\lvert\arcPosi\setminus\wrapTc\rvert = \lvert(\wrapSet\setminus\wrapTc)\cup\arcNega\rvert$

Therefore, we have:
\[
\begin{array}{rl}
\because & \wrapTc = \arcPosi \cap \wrapSet,\; \wrapSet \cap \arcNega = \emptyset \\ [4pt]
\therefore & \wrapTc \cap \arcNega = \emptyset \\ [4pt]
\therefore & (\wrapSet \setminus \wrapTc) \cup \arcNega = \emptyset \\ [4pt]
\therefore & \lvert\wrapSet\setminus\wrapTc\cup\arcNega\rvert = \lvert\wrapSet\setminus\wrapTc\biguplus\arcNega\rvert = \lvert\wrapSet\setminus\wrapTc\rvert + \lvert\arcNega\rvert \\ [4pt]
\because & \lvert\arcPosi\setminus\wrapTc\rvert = \lvert(\wrapSet\setminus\wrapTc)\cup\arcNega\rvert \\ [4pt]
\therefore & \lvert\wrapSet\setminus\wrapTc\rvert + \lvert\arcNega\rvert = \lvert\arcPosi\setminus\wrapTc\rvert \\ [4pt]
\because & \begin{cases}
    \wrapSet = \wrapTc \biguplus (\wrapSet\setminus\wrapTc) \implies \lvert\wrapSet\rvert = \lvert\wrapTc\rvert + \lvert\wrapSet\setminus\wrapTc\rvert \\
    \arcPosi = \wrapTc \biguplus (\arcPosi\setminus\wrapTc) \implies \lvert\arcPosi\rvert = \lvert\wrapTc\rvert + \lvert\arcPosi\setminus\wrapTc\rvert
\end{cases} \\ [14pt]
\therefore & (\lvert\wrapSet\setminus\wrapTc\rvert + \lvert\arcNega\rvert) + \lvert\wrapTc\rvert = (\lvert\arcPosi\setminus\wrapTc\rvert) + \lvert\wrapTc\rvert \quad\text{(add $\lvert\arcNega\rvert$ in both sides)}\\ [4pt]
\implies & \lvert\wrapSet\rvert + \lvert\arcNega\rvert = \lvert\arcPosi\rvert \\ [4pt]
\implies & \lvert\arcPosi\rvert = \lvert\arcNega\rvert + \wrapnum_{fk} \\ [4pt]
\end{array}
\]
The conclusion follows

% Traverse $\widetilde{\pathInRoute{f}{k}}$ once in its chronological order and record the binary sequence
% \[
%     \sign_r :=
%     \begin{cases}
%         1, & t_o(a_{k,r})\succ_{t_o(a_{k,r})}\checkTime,\\
%         0, & \text{otherwise},
%     \end{cases}
% \qquad r=1,\ldots,Q_k.
% \]
% Because consecutive arcs share endpoints, $t_o(a_{k,r\oplus1})=t_d(a_{k,r})$ (with indices modulo $Q_k$), hence a change in $\sign_r$ occurs exactly when the interval traced by $a_{k,r}$ crosses $t_c$ in the forward direction of time measured from $t_o(a_{k,r})$.
% If $a_{k,r}\in\arcPosi\setminus\wrapSet_{fk}$ then $t_o(a_{k,r})\le t_c < t_d(a_{k,r})$ and the arc does not wrap around the horizon, which implies $\sign_r=0$ and $\sign_{r\oplus1}=1$.
% If $a_{k,r}\in\arcNega$ then $t_d(a_{k,r})\le t_c < t_o(a_{k,r})$ so $\sign_r=1$ and $\sign_{r\oplus1}=0$.
% In every remaining case--including wrap-around arcs in $\wrapSet_{fk}$--the cut is not crossed in this sense and $\sign_r=\sign_{r\oplus1}$.
% Consequently, for each $r$ we have
% \[
%     \sign_{r\oplus1}-\sign_r
%     = \begin{cases}
%         +1, & a_{k,r}\in\arcPosi\setminus\wrapSet_{fk},\\
%         -1, & a_{k,r}\in\arcNega,\\
%         0,  & \text{otherwise}.
%     \end{cases}
% \]
% Summing these identities over $r=1,\ldots,Q_k$ gives
% $0 = |\arcPosi\setminus\wrapSet_{fk}|-|\arcNega|$, so that $|\arcPosi|=|\arcNega|+|\wrapSet_{fk}|$.
% The conclusion follows because $|\wrapSet_{fk}|=\wrapnum_{fk}$ by definition.
\end{proof}